\documentclass[11pt,a4paper]{article}
\usepackage[margin=2.4cm]{geometry}
\usepackage[T1]{fontenc}
\usepackage{lmodern}
\usepackage[hidelinks]{hyperref}
\usepackage{parskip}
\pagestyle{empty}

\begin{document}

Nano4Imaging GmbH\\
Attn: Samuel Mueller, Recruiting Team\\
Merowingerplatz 1\\
Life Science Center\\
40225 Duesseldorf\\
Germany

\begin{flushright}
17 August 2026
\end{flushright}

Dear Mr Mueller,

I am applying for the AI Research Engineer position because it brings together the two parts of my work that I care about most: developing strong computer-vision methods and making them useful in real clinical settings. My background is in machine learning for medical imaging, and I have worked with clinicians and research partners from my first substantial projects through to systems that are now used in clinical trials.

My interest in this area started during my studies, when I worked on projects in cooperation with a local hospital. Medical applications have stayed especially motivating to me because the work can have a direct impact on how patients are diagnosed and treated, and because the results are shaped by close cooperation with the people who care for them. I continued in this direction with my thesis on anatomical landmark localization for orthopedics. In cooperation with spine specialists, I improved the previous model by 30\%; the work was presented at DWG 2023.

Since then, I have continued to work across different medical-imaging problems. In the three-year BMBF Photiomics project, I developed 3D instance-segmentation and registration models for liver-cancer research, successfully carrying the project through despite not having a background in cellular science. More recently, I developed, validated, and deployed an end-to-end 2D--3D registration system for implant pose estimation from X-rays. The work was presented at EFORT 2026, and the system is now used at scale in current clinical trials.

This experience is a good match for Nano4Imaging's work on real-time MRI navigation. I have wanted to gain hands-on experience with MRI for a long time and am excited to take on the challenges of this modality. My work with anisotropic 3D light-sheet microscopy data has given me practical experience with voxel-based volumetric data, differences in spatial resolution, and 3D image processing. I believe that foundation would transfer well to MRI while still leaving room to learn the modality-specific physics and artifacts. I would be excited to bring this background, together with my experience in deep learning, registration, segmentation, and clinical collaboration, to the development of TRACKR.

Thank you for considering my application. I would welcome the opportunity to discuss how my experience could contribute to Nano4Imaging's work on safer image-guided interventions.

Kind regards,\\
Kostiantyn Pysanyi

\end{document}
